% Source: https://artofproblemsolving.com/community/c3h2097799_help_me_with_this_sequence_i_cant_figure_it_out
% Stable identity: aops-topic-2097799-q01
\textbf{(a)} From the first three terms, $8=3p+q$ and $23=8p+q$. Subtracting gives $15=5p$, so $p=3$ and $q=-1$. The rule $a_{n+1}=3a_n-1$ also fits the remaining terms: $3\cdot23-1=68$ and $3\cdot68-1=203$. The next term is $3\cdot203-1=\boxed{608}$.

\textbf{(b)} The fixed point of $t\mapsto3t-1$ is $t=\dfrac{1}{2}$, and
\[a_{n+1}-\dfrac{1}{2}=3a_n-1-\dfrac{1}{2}=3\left(a_n-\dfrac{1}{2}\right).\]
So $a_n-\dfrac{1}{2}$ is a geometric sequence with ratio $3$ and first term $3-\dfrac{1}{2}=\dfrac{5}{2}$, and
\[a_n=\boxed{\dfrac{5\cdot3^{\,n-1}+1}{2}}.\]
Check: $n=1,2,3$ give $3$, $8$, $23$.
